\bookchapter{第五十二章：4大コンポーネントとActivityのライフサイクル}

\chapterintro{Composeをアプリの「顔」とするなら、この章で扱うのはアプリの「骨格」です。Androidシステムには4つの基本コンポーネントがあり、すべてのアプリはそれらの組み合わせでできています。中でもいちばん大事なActivityには、すべてのAndroid開発者が脳に刻むべき「ライフサイクル図」があります。この章を終えれば、「バックグラウンドに行って戻ったら、なぜアプリの状態が消えているのか」といった永遠の謎が解けるようになります。}

\booksection{4大コンポーネント：Androidの「大黒柱」}

Androidは公式に4種類の基本コンポーネントを定めていて、それぞれが1種類の仕事を担います。

\begin{center}
\begin{tabularx}{\linewidth}{|X|X|X|}
\hline
\textbf{コンポーネント} & \textbf{役割} & \textbf{たとえるなら} \tabularnewline
\hline
Activity & 目に見える画面 & お店の店構え——ユーザーが見て、操作できる部分 \tabularnewline
Service & バックグラウンドで働く「見えない人」 & 倉庫の作業員——音楽を流す、ファイルをダウンロードする。画面は不要 \tabularnewline
BroadcastReceiver & システムのブロードキャストを受け取る & 受付——「電池が残り少ない」「ネットが切れた」を告げるスピーカー \tabularnewline
ContentProvider & アプリをまたいでデータを共有 & 共用の資料棚——たとえば連絡先を別のアプリに共有する \tabularnewline
\hline
\end{tabularx}
\end{center}

\begin{center}\small\color{gray}Android 4大コンポーネント早見\end{center}

初心者のうちは、時間の95\%をActivityとの付き合いに使います。ServiceとBroadcastReceiverは音楽再生やプッシュ通知で出番があり、ContentProviderはいちばん登場が少ないです（よく使うデータはシステムがすでにまとめてくれています）。もう1つ大事なこと。4大コンポーネントは、すべてAndroidManifest.xmlに「住民登録」が要ります。登録しないと、システムはそれを見つけられません。

\booksection{Activityのライフサイクル：Android開発で最初の「宝の地図」}

スマホのアプリは、いつでも中断され得ます。電話がかかってくる、ユーザーがWeChatの返信に切り替える、メモリが足りなくてシステムに止められる……Activityの一生は、この図の上をぐるぐる巡る旅なのです。

onCreate（誕生：初期化）

onStart（見えるようになった）

onResume（前面で操作中）

onPause（一部が隠れた）

onStop（まったく見えない）

onDestroy（破棄）

onRestart / onStart / onResume：前面に戻るときは同じ道をたどる

\begin{center}\small\color{gray}図4：Activityのライフサイクル（進む向き：上段が誕生 下段が退場。戻るときは同じ道を逆にたどる）\end{center}

\begin{center}
\begin{tabularx}{\linewidth}{|X|X|X|}
\hline
\textbf{コールバック} & \textbf{いつ呼ばれる？} & \textbf{典型的な用途} \tabularnewline
\hline
onCreate & Activityが生成されるとき（1回だけ） & UIの読み込み、データの初期化 \tabularnewline
onStart & 画面が見える状態になる & 軽めの処理を始める \tabularnewline
onResume & フォーカスを得て、操作できるようになる & アニメーションの開始、再生の再開 \tabularnewline
onPause & 一部が隠れる／切り替わる瞬間 & アニメーションの一時停止、下書きの保存 \tabularnewline
onStop & 完全に見えなくなる & 重いリソースの解放、ダウンロードの一時停止 \tabularnewline
onDestroy & 破棄される前 & 最後のクリーンアップ \tabularnewline
\hline
\end{tabularx}
\end{center}

\begin{center}\small\color{gray}6つのライフサイクルコールバックの典型的な用途\end{center}

\textbf{落とし穴注意：画面回転はActivityを「殺して」しまう}\par

初心者にいちばん多い失敗例：アプリの画面を回したら、入力欄の文字が全部消えた。原因：画面を回転すると、デフォルトではActivityが破棄されて再び作り直されるためです（onDestroy onCreate が一巡します）。解決策：rememberSaveable で状態を保存する（Composeでのrememberの代わり）。再生成のときに自動で値を取り戻してくれます。第6·7章で何度も使うことになります。

\booksection{Intent：アプリの内と外をつなぐ「切符」}

Intent（インテント）は、コンポーネント同士がやりとりするための「切符」です。ある画面から別の画面へ移動したいときも、ブラウザを開きたいときも、カメラを起動したいときも、必ず「切符を買って乗車」します。

\begin{lstlisting}[language=Kotlin]
// 方法その1：明示的Intent——行き先のページをはっきり指定
val intent = Intent(this, DetailActivity::class.java)
intent.putExtra("name", "たろう")     // 車に「荷物」を積むイメージ（データ）
startActivity(intent)

// 方法その2：暗黙的Intent——何をしたいかだけ言う。実行役はシステムが探す
val web = Intent(Intent.ACTION_VIEW, Uri.parse("https://developer.android.com"))
startActivity(web)   // システムがブラウザを開く/または選択させる
\end{lstlisting}

2つの方式にはそれぞれ役割があります。明示的Intentは自分のアプリ内の画面遷移に、暗黙的Intentはシステムや他アプリの機能を呼び出すとき（Webページを開く、電話をかける、共有するなど）に使います。Composeプロジェクトでは、Composableの中で LocalContext.current でContextを取れば startActivity が使えます。

\booksection{複数画面とナビゲーション}

実際のアプリには複数の画面があります。Composeの公式が推奨するナビゲーション手法はNavigation Compose（依存ライブラリ androidx.navigation）です。

\begin{lstlisting}[language=Kotlin]
val navController = rememberNavController()

NavHost(navController, startDestination = "home") {
    composable("home") { HomeScreen(navController) }
    composable("detail/{name}") { entry ->
        val name = entry.arguments?.getString("name")
        DetailScreen(name ?: "不明")
    }
}

// ページ内で遷移：
navController.navigate("detail/たろう")     // 引数付き
navController.popBackStack()              // 前のページへ戻る
\end{lstlisting}

\begin{tcbinfo}{ワンポイント：画面が増えてきたら？}
プロジェクトが大きくなったら、すべての画面のルートを1つのファイル（たとえば AppNavGraph.kt）にまとめて一元管理しましょう。「どこへ遷移するか」というロジックが各画面に散らばらないようにするのです。これが「動くもの」から「保守できるもの」へ進む第一歩です。

\end{tcbinfo}

\booksection{ついでに知っておく：Fragmentとは何か（そして本書で扱わない理由）}

古いチュートリアルを読むと必ずFragmentに出会います。FragmentはActivityの中に置ける「再利用できる画面の部品」で、昔ながらの複数画面アーキテクチャでは主役でした。しかしComposeの時代には、画面はComposable関数＋Navigationのルートで足りていて、公式もFragmentの立ち位置を弱めています。そのため本書では本筋に含めていません。

アドバイス：新しいプロジェクトでFragmentを進んで使う必要はありません。ただし古いコードを読むときや、面接で聞かれたときには「Activityに組み込む画面の部品」だと分かっていれば十分です。学習リソースは常に2023年以降のものを優先しましょう。Androidの技術スタックの「古い常識」は、想像以上に速く役目を終えていきます。

\begin{tcbpractice}{やってみよう：カウンターに「記憶」をつける}
第五十一章のカウンターを改造してみましょう。(1) remember を rememberSaveable に替え、画面を回転しても数字が残ることを確認する。(2)「このアプリについて」ボタンを追加し、Navigationで2つ目の画面へ遷移してバージョン番号を表示する。(3) 両方の画面の onCreate/onPause にログ出力（Log.d）を入れ、Logcatでライフサイクルが呼ばれる順番を観察する。この3つができれば、ライフサイクルはもう机上の知識ではありません。

\end{tcbpractice}

\begin{tcbsummary}{本章のまとめ}
\begin{itemize}
  \item 4大コンポーネント：Activity（顔）·Service（バックグラウンド）·BroadcastReceiver（スピーカー）·ContentProvider（資料棚）
  \item ライフサイクルの6つのコールバック：onCreate onStart onResume onPause onStop onDestroy
  \item 画面回転でActivityは破棄·再生成される：rememberSaveableで状態を守る
  \item Intentは「切符」：明示的は自アプリの画面遷移、暗黙的はシステム機能の呼び出し
  \item 複数画面はNavigation Composeで：NavHost＋composableルート
\end{itemize}

\end{tcbsummary}
